On veut préparer un volume $V=\qty{100}{\mL}$ d'une solution de sulfate
d'aluminium anhydre (\ch{Al2(SO4)3}) de concentration $c=\qty{0.050}{\mol\per\L}$.

Pour cela on pèse une masse $m$ de sulfate d'aluminium anhydre
que l' on introduit dans une fiole jaugée et que l'on complète avec de l'eau
distillée jusqu'au trait de jauge.
\begin{enumerate}
    \item Écrire l'équation de dissolution du soluté dans l'eau.~\textbf{(1pt)}
    \item Pourquoi chaque ion s'entoure t-il de molécule d'eau? Comment appelle
          t-on ce phénomène ? Prenez l'exemple de l'ion \ch{Al^{3+}_{(aq)}} et
          dessinez les molécules d'eau autour de cet ion.~\textbf{(2pt)}
    \item Calculer la concentration molaire des ions \ch{Al^{3+}_{(aq)}} et
          \ch{SO4^{2-}_{(aq)}}.~\textbf{(1pt)}
    \item Exprimer la concentration molaire $c$ du soluté (sulfate d'aluminium)
          en fonction de $m$, $V$ et $M(\ch{Al2(SO4)3})$.~\textbf{(1pt)}
    \item Calculer la masse $m$ de soluté à peser.~\textbf{(0,5pt)}
\end{enumerate}

\begin{donnees}
    en $\unit{\g\per\mol}$~: $M(\ch{Al})=27,0$~; $M(\ch{S})=32,1$~;
    $M(\ch{O})=16,0$.
\end{donnees}
